% Intermediate summary matrix for Q3. May later be condensed, moved to an appendix, or redrawn as a figure.
\begin{table}[tbp]
\centering
\footnotesize
\singlespacing
\caption{Working summary for Q3: what neural measures can and cannot show. Each row separates a warranted inference from a common over-reading.}
\label{tab:summary-q3}
\begin{tabular}{@{}p{1.45in}p{2.4in}p{2.4in}@{}}
\toprule
\textbf{Observation} & \textbf{Warranted inference} & \textbf{Not warranted} \\
\midrule
Amygdala differs for 22- vs 50-kHz calls; matched tones reproduce most of it & Acoustic detection / frequency code & Call-specific meaning; threat encoding \\
\addlinespace[2pt]
50-kHz raises, 22-kHz does not raise, NAc dopamine & Valence-sensitive transient & Unitary reward or fear signal \\
\addlinespace[2pt]
Noise lowers NAc DA, raises tail-striatum DA; intensity scaling depends on preparation & Region- and context-specific responses & Scalar threat code in NAc \\
\addlinespace[2pt]
Chirps reduce c-Fos in medial amygdala and superior colliculus & Threat $\neq$ more activation & Mapping to field escape/vigilance \\
\addlinespace[2pt]
IEG patterns differ by call type & Regional engagement over minutes & Detection vs interpretation vs action \\
\addlinespace[2pt]
DA to ecological threat calls & Sparse / largely untested & --- \\
\bottomrule
\end{tabular}
\end{table}
